%==============================================================================
% Write-Up de CTF - Crack The Hash (TryHackMe)
% Materia : CIC0087 - Topicos Avancados em Computadores - Turma 04 - 2026/2
% Depto.  : Departamento de Ciencias da Computacao
% Baseado na "Metodologia para Construcao de Write-Ups em CTF"
%
% Compilacao:  pdflatex main.tex   (rodar 2x para atualizar o Sumario)
%==============================================================================
\documentclass[12pt,a4paper]{article}

%------------------------- Codificacao e idioma -------------------------------
\usepackage[utf8]{inputenc}
\usepackage[T1]{fontenc}
\usepackage[brazil]{babel}

%------------------------- Fonte Times New Roman ------------------------------
\usepackage{newtxtext}   % corpo em Times
\usepackage{newtxmath}   % matematica em Times

%------------------------- Margens (minimo 2 cm) ------------------------------
\usepackage[a4paper,margin=2cm,headheight=15pt,includehead,includefoot]{geometry}

%------------------------- Utilitarios ----------------------------------------
\usepackage{graphicx}
\graphicspath{{img/}{./img/}}
\usepackage{xcolor}
\usepackage{float}
\usepackage[hidelinks]{hyperref}
\usepackage{fancyhdr}
\usepackage{listings}
\usepackage[most]{tcolorbox}
\usepackage{caption}
\usepackage{enumitem}
\usepackage{setspace}
\usepackage{array}
\usepackage{longtable}

\hypersetup{
  colorlinks=true,
  linkcolor=black,
  urlcolor=blue!60!black,
  citecolor=black
}
\onehalfspacing

%============================= METADADOS =======================================
\newcommand{\desafioNome}{Crack The Hash}
\newcommand{\plataforma}{TryHackMe}
\newcommand{\dificuldade}{Fácil}
\newcommand{\dataResolucao}{10/09/2026}
\newcommand{\autorNome}{Luís Eduardo Curi Serra}

% Dados academicos (fixos para a disciplina)
\newcommand{\materia}{CIC0087 -- Tópicos Avançados em Computadores -- Turma 04 -- 2026/2}
\newcommand{\departamento}{Departamento de Ciências da Computação}
\newcommand{\matricula}{251033574}
\newcommand{\professor}{Roberto Rodrigues Filho}
\newcommand{\universidade}{Universidade de Brasília}
%===============================================================================

%------------------------- Cabecalho e rodape ---------------------------------
\pagestyle{fancy}
\fancyhf{}
\lhead{\small\itshape \desafioNome}
\rhead{\small\itshape \autorNome}
\lfoot{\small \materia}
\rfoot{\small Página \thepage}
\renewcommand{\headrulewidth}{0.4pt}
\renewcommand{\footrulewidth}{0.4pt}

%------------------------- Estilo dos blocos de codigo ------------------------
\definecolor{codebg}{gray}{0.95}
\definecolor{codeframe}{gray}{0.75}
\definecolor{cmt}{RGB}{100,120,100}
\definecolor{kw}{RGB}{0,0,150}
\definecolor{str}{RGB}{150,50,50}

\lstdefinestyle{shell}{
  backgroundcolor=\color{codebg},
  basicstyle=\ttfamily\small,
  keywordstyle=\color{kw}\bfseries,
  commentstyle=\color{cmt}\itshape,
  stringstyle=\color{str},
  frame=single,
  rulecolor=\color{codeframe},
  breaklines=true,
  breakatwhitespace=false,
  postbreak=\mbox{\textcolor{codeframe}{$\hookrightarrow$}\space},
  columns=fullflexible,
  keepspaces=true,
  showstringspaces=false,
  tabsize=2,
  captionpos=b,
  xleftmargin=0.4em,
  xrightmargin=0.4em,
  aboveskip=1em,
  belowskip=1em,
}
\lstset{style=shell}

% Atalho: comando de terminal em linha
\newcommand{\cmd}[1]{\lstinline[style=shell]|#1|}

% Ambiente para OUTPUT (fundo levemente diferente, fonte menor)
\lstdefinestyle{output}{
  style=shell,
  backgroundcolor=\color{gray!8},
  basicstyle=\ttfamily\footnotesize,
  keywordstyle=\ttfamily,
}

%------------ Caixa de "becos sem saida" (tentativas sem sucesso) --------------
\definecolor{deadendbg}{RGB}{255,248,240}
\definecolor{deadendframe}{RGB}{193,108,44}
\newtcolorbox{becosemsaida}[1][Tentativas sem sucesso (becos sem saída)]{%
  enhanced, breakable,
  colback=deadendbg, colframe=deadendframe,
  colbacktitle=deadendframe, coltitle=white,
  boxrule=0.6pt, arc=2pt,
  left=8pt, right=8pt, top=6pt, bottom=6pt,
  fonttitle=\bfseries\small,
  title={#1}%
}

%===============================================================================
\begin{document}

%============================= 2.1 CAPA =========================================
\begin{titlepage}
  \thispagestyle{empty}
  \centering
  {\scshape \universidade \par}
  {\scshape \departamento \par}
  \vspace{0.4cm}
  {\small \materia \par}
  \vspace{2.5cm}

  \rule{\linewidth}{0.4pt}\par
  \vspace{0.6cm}
  {\huge \bfseries Write-Up: \desafioNome \par}
  \vspace{0.4cm}
  {\large Segurança Ofensiva --- Capture The Flag \par}
  \vspace{0.6cm}
  \rule{\linewidth}{0.4pt}\par

  \vspace{2.0cm}
  \begin{tabular}{rl}
    \textbf{Plataforma:}       & \plataforma       \\[4pt]
    \textbf{Dificuldade:}      & \dificuldade      \\[4pt]
    \textbf{Data de resolução:}& \dataResolucao    \\
  \end{tabular}

  \vfill

  \begin{tabular}{rl}
    \textbf{Autor:}     & \autorNome  \\[4pt]
    \textbf{Matrícula:} & \matricula  \\[4pt]
    \textbf{Professor:} & \professor  \\
  \end{tabular}

  \vspace{1.2cm}
  {\small \today \par}
\end{titlepage}

%============================= 2.2 SUMARIO ======================================
\newpage
\tableofcontents
\thispagestyle{fancy}
\newpage

%============================= 2.3 VISAO GERAL ==================================
\section{Visão Geral}
\textbf{Crack The Hash} é uma sala de \emph{criptografia} do TryHackMe, de
dificuldade \emph{fácil}. Diferentemente de uma máquina, não há serviço a
explorar nem escalada de privilégio: o desafio entrega uma série de \emph{hashes}
e o objetivo é \textbf{identificar o algoritmo} de cada um e \textbf{recuperar o
texto-claro}. A sala está dividida em dois níveis de dificuldade crescente, além
disso foram resolvidos \textbf{hashes adicionais} propostos como exercício complementar. A Tabela~\ref{tab:resumo}

O fluxo aplicado a cada hash é sempre o mesmo: (i) \textbf{identificar} o
algoritmo pelo comprimento e formato do hash, com apoio de \texttt{hashid},
\texttt{hash-identifier} e do \emph{Magic} do CyberChef; (ii) \textbf{mapear o
modo} correspondente do \texttt{hashcat} (\cmd{-m}); e (iii) \textbf{atacar} com
o ataque adequado (\cmd{-a}) --- dicionário (\texttt{rockyou.txt}), dicionário
com regras de mutação, ou máscara/força-bruta quando há dica sobre o formato da
senha. Como os hashes são a entrada pública do desafio e as senhas recuperadas
são as próprias respostas didáticas da sala, ambos são exibidos em claro neste
documento.

%==================== 2.4 RECONHECIMENTO / FERRAMENTAS =========================
\section{Reconhecimento: Identificação e Ferramentas}

O ``reconhecimento'' de um desafio de \emph{cracking} consiste em determinar,
para cada hash, qual algoritmo o gerou.
O comprimento em caracteres hexadecimais é a primeira pista (32 $\rightarrow$
família MD5/MD4; 40 $\rightarrow$ SHA-1; 64 $\rightarrow$ SHA-256); prefixos como
\cmd{$2y$} (bcrypt) e \cmd{$6$} (sha512crypt) são identificadores diretos do
formato.

\begin{itemize}[leftmargin=1.4em,itemsep=2pt,topsep=2pt]
  \item \textbf{hashid} --- lista os candidatos de algoritmo e já informa o modo
        (\cmd{-m}) do \texttt{hashcat}: \cmd{hashid -m '<hash>'}.
  \item \textbf{hash-identifier} --- ferramenta interativa equivalente.
  \item \textbf{CyberChef (\emph{Magic})} --- detecta o tipo e sugere a família
        de algoritmos a partir do formato.
\end{itemize}

Um ponto essencial é que ferramentas de identificação retornam \emph{várias}
possibilidades para um mesmo comprimento (p.\,ex., um hash de 32 hex pode ser
MD5, MD4, NTLM, RIPEMD-128 etc.). O procedimento correto é testar os candidatos
em ordem de probabilidade até um deles quebrar --- foi exatamente o que
distinguiu os casos mais trabalhosos (Nível 1 \#5 e Nível 2 \#2).

\subsection{Ferramenta de quebra --- hashcat}
A quebra propriamente dita foi feita com \texttt{hashcat}. Os dois parâmetros
centrais são o \textbf{modo} (\cmd{-m}, o algoritmo) e o \textbf{tipo de ataque}
(\cmd{-a}). A Tabela~\ref{tab:modos} reúne os modos efetivamente usados no
desafio.

\begin{table}[H]
  \centering
  \caption{Modos de \texttt{hashcat} utilizados neste desafio.}
  \label{tab:modos}
  \begin{tabular}{>{\ttfamily}c l l}
    \hline
    \textbf{-m} & \textbf{Algoritmo} & \textbf{Onde apareceu} \\
    \hline
    0    & MD5                         & N1\#1, Add\#1 \\
    100  & SHA-1                       & N1\#2, Add\#2 \\
    900  & MD4                         & N1\#5, Add\#3 \\
    1000 & NTLM                        & N2\#2 \\
    1400 & SHA-256                     & N1\#3, N2\#1, Add\#5 \\
    1800 & sha512crypt (\$6\$)         & N2\#3 \\
    3200 & bcrypt (\$2*\$)             & N1\#4 \\
    160  & HMAC-SHA1 (chave = sal)     & N2\#4, Add\#4 \\
    \hline
  \end{tabular}
\end{table}

\noindent Tipos de ataque (\cmd{-a}) empregados: \textbf{0} (dicionário),
\textbf{3} (máscara / força-bruta). Recursos auxiliares:

\begin{itemize}[leftmargin=1.4em,itemsep=2pt,topsep=2pt]
  \item \textbf{Dicionário:} \cmd{/usr/share/wordlists/rockyou.txt}.
  \item \textbf{Regras} (\cmd{-r}): geram variações de cada palavra do dicionário
        (maiúsculas, números no fim, \emph{leetspeak}), usadas quando a senha é
        uma palavra comum mutada. Empregou-se \cmd{best66.rule}, nativo do \emph{Kali Linux}.
  \item \textbf{Máscaras} (\cmd{-a 3}): \cmd{?d}=dígitos, \cmd{?l}=minúsculas,
        \cmd{?u}=maiúsculas. Com \cmd{-1} define-se um \emph{charset} próprio
        referenciado por \cmd{?1} --- útil quando cada posição pode ser de vários
        conjuntos.
  \item \textbf{Sal:} para hashes com sal, o \texttt{hashcat} recebe o par no
        formato \cmd{hash:sal} (modos 110/120/160), ou o sal já vem embutido no
        próprio hash (formato \cmd{$6$<sal>$<hash>}).
  \item \textbf{Otimização} (\cmd{-O}): habilita o \emph{kernel} otimizado
        (limita o tamanho da senha, mas acelera bastante).
\end{itemize}

%============================= 2.5 EXPLORACAO ===================================
\section{Exploração: Quebra dos Hashes}

Para cada hash apresenta-se: o valor de entrada, o algoritmo identificado e o
respectivo modo, o comando de quebra e o texto-claro recuperado. Onde houve
tentativas frustradas relevantes, elas são registradas em destaque, pois
compõem o raciocínio de identificação.

%------------------------------- NIVEL 1 --------------------------------------
\subsection{Nível 1 --- (Task 1)}

\begin{table}[H]
  \centering
  \caption{Resumo do Nível 1.}
  \label{tab:n1}
  \begin{tabular}{l l >{\ttfamily}c l l}
    \hline
    \textbf{Hash (abrev.)} & \textbf{Algoritmo} & \textbf{-m} & \textbf{Ataque} & \textbf{Texto-claro} \\
    \hline
    48bb\ldots ed4d   & MD5     & 0    & dicionário            & \texttt{easy} \\
    CBFD\ldots 2A97   & SHA-1   & 100  & dicionário            & \texttt{password123} \\
    1C8B\ldots 3032   & SHA-256 & 1400 & dicionário            & \texttt{letmein} \\
    \$2y\$12\$\ldots  & bcrypt  & 3200 & dicionário (4 letras) & \texttt{bleh} \\
    2794\ldots 3daa   & MD4     & 900  & dicionário + regras   & \texttt{Eternity22} \\
    \hline
  \end{tabular}
\end{table}

\subsubsection{\#1 --- MD5}
Hash de 32 hex: \cmd{48bb6e862e54f2a795ffc4e541caed4d}. O CyberChef aponta a
família MD (MD5, MD4, RIPEMD-128\ldots); começa-se pelo mais provável, MD5
(\cmd{-m 0}), com dicionário.

\begin{lstlisting}[language=bash,caption={Quebra do hash MD5 por dicionário.}]
hashcat -m 0 -a 0 t1_1.txt /usr/share/wordlists/rockyou.txt
\end{lstlisting}
\noindent Texto-claro: \textbf{easy}.

\subsubsection{\#2 --- SHA-1}
Hash de 40 hex: \cmd{CBFDAC6008F9CAB4083784CBD1874F76618D2A97}. O comprimento
identifica SHA-1 (\cmd{-m 100}).

\begin{lstlisting}[language=bash,caption={Quebra do hash SHA-1 por dicionário.}]
hashcat -m 100 -a 0 t1_2.txt /usr/share/wordlists/rockyou.txt
\end{lstlisting}
\noindent Texto-claro: \textbf{password123}.

\subsubsection{\#3 --- SHA-256}
Hash de 64 hex: \cmd{1C8BFE8F801D79745C4631D09FFF36C82AA37FC4CCE4FC946683D7B336B63032}.
Comprimento típico de SHA-256 (\cmd{-m 1400}).

\begin{lstlisting}[language=bash,caption={Quebra do hash SHA-256 por dicionário.}]
hashcat -m 1400 -a 0 t1_3.txt /usr/share/wordlists/rockyou.txt
\end{lstlisting}
\noindent Texto-claro: \textbf{letmein}.

\subsubsection{\#4 --- bcrypt}
Hash \cmd{$2y$12$Dwt1BZj6pcyc3Dy1FWZ5ieeUznr71EeNkJkUlypTsgbX1H68wsRom}. O
prefixo \cmd{$2y$} e o fator de custo \cmd{12} identificam \texttt{bcrypt}
(\cmd{-m 3200}) via \texttt{hashid}. Por ser um algoritmo deliberadamente lento
($2^{12}$ iterações por candidato), o dicionário completo é inviável em tempo
hábil; por isso o \texttt{rockyou} foi filtrado para conter apenas palavras
alfabéticas de \textbf{4 caracteres}, reduzindo drasticamente o espaço de busca.

\begin{lstlisting}[language=bash,caption={Filtragem do dicionário e quebra do bcrypt.}]
grep -xE '[A-Za-z]{4}' /usr/share/wordlists/rockyou.txt > words4.txt
hashcat -m 3200 -a 0 t1_4.txt words4.txt
\end{lstlisting}
\noindent Texto-claro: \textbf{bleh}.

\subsubsection{\#5 --- MD4}
Hash de 32 hex: \cmd{279412f945939ba78ce0758d3fd83daa}. Foi o hash mais
trabalhoso do Nível 1: o texto-claro (\textbf{Eternity22}) não está no
\texttt{rockyou} puro e o algoritmo não é MD5. Após descartar MD5 e SHA-1,
confirma-se \textbf{MD4} (\cmd{-m 900}); a senha é uma palavra do dicionário
\emph{mutada}, recuperada com o dicionário mais o conjunto de regras
\cmd{best66.rule}.

\begin{lstlisting}[language=bash,caption={Quebra do hash MD4 com dicionário + regras.}]
hashcat -m 900 -a 0 t1_5.txt /usr/share/wordlists/rockyou.txt \
        -r /usr/share/hashcat/rules/best66.rule
\end{lstlisting}
\noindent Texto-claro: \textbf{Eternity22}.

\begin{becosemsaida}
\begin{itemize}[leftmargin=1.4em,itemsep=2pt,topsep=2pt]
  \item MD5 (\cmd{-m 0}) e MD4 (\cmd{-m 900}) com \texttt{rockyou} \emph{puro}:
        sem resultado --- a senha é uma mutação, não uma palavra literal.
  \item Máscara de força-bruta com 10 minúsculas
        (\cmd{-a 3 ?l?l?l?l?l?l?l?l?l?l}): estimativa de $\approx 5$ dias mesmo
        com \cmd{-O} --- inviável, abandonado (e a senha nem é toda minúscula).
  \item Lista externa \emph{fortinet-2021} (SecLists): também sem resultado.
  \item Só a combinação \textbf{MD4 + \texttt{rockyou} + regras} destravou.
\end{itemize}
\end{becosemsaida}

%------------------------------- NIVEL 2 --------------------------------------
\subsection{Nível 2 --- (Task 2)}

\begin{table}[H]
  \centering
  \caption{Resumo do Nível 2.}
  \label{tab:n2}
  \begin{tabular}{l l >{\ttfamily}c l l}
    \hline
    \textbf{Hash (abrev.)} & \textbf{Algoritmo} & \textbf{-m} & \textbf{Ataque} & \textbf{Texto-claro} \\
    \hline
    F09E\ldots 0C85          & SHA-256      & 1400 & dicionário + regras   & \texttt{paule} \\
    1DFE\ldots CC1B          & NTLM         & 1000 & dicionário + regras   & \texttt{n63umy8lkf4i} \\
    \$6\$\ldots As02.        & sha512crypt  & 1800 & dicionário (6 chars)  & \texttt{waka99} \\
    e5d8\ldots 56d6 : sal    & HMAC-SHA1    & 160  & dicionário + regras   & \texttt{481616481616} \\
    \hline
  \end{tabular}
\end{table}

\subsubsection{\#1 --- SHA-256 (com regras)}
Hash de 64 hex: \cmd{F09EDCB1FCEFC6DFB23DC3505A882655FF77375ED8AA2D1C13F640FCCC2D0C85}
(SHA-256, \cmd{-m 1400}). O \texttt{rockyou} puro não encontra; a senha é uma
mutação, quebrada com regras.

\begin{lstlisting}[language=bash,caption={SHA-256 com dicionário + regras.}]
hashcat -m 1400 -a 0 t2_1.txt /usr/share/wordlists/rockyou.txt \
        -r /usr/share/hashcat/rules/best66.rule
\end{lstlisting}
\noindent Texto-claro: \textbf{paule}.

\subsubsection{\#2 --- NTLM}
Hash de 32 hex: \cmd{1DFECA0C002AE40B8619ECF94819CC1B}. O \texttt{hashid} retorna
uma lista extensa (MD5, MD4, \emph{Double MD5}, LM, NTLM\ldots). Testando os
candidatos em ordem, todos falham até \textbf{NTLM} (\cmd{-m 1000}), que quebra.
NTLM é o hash de senhas do Windows.

\begin{lstlisting}[language=bash,caption={NTLM com dicionário + regras.}]
hashcat -m 1000 -a 0 t2_2.txt /usr/share/wordlists/rockyou.txt \
        -r /usr/share/hashcat/rules/best66.rule
\end{lstlisting}
\noindent Texto-claro: \textbf{n63umy8lkf4i}.

\begin{becosemsaida}
Candidatos testados antes do NTLM, todos sem sucesso:
\begin{itemize}[leftmargin=1.4em,itemsep=2pt,topsep=2pt]
  \item MD5 (\cmd{-m 0}) e MD4 (\cmd{-m 900}) --- com \cmd{-O} e regras.
  \item \emph{Double MD5} (\cmd{-m 2600}) e LM (\cmd{-m 3000}).
  \item \emph{Lotus Notes/Domino 5} (\cmd{-m 8600}): estimativa de $\approx 30$
        min --- descartado por custo/benefício.
\end{itemize}
\end{becosemsaida}

\subsubsection{\#3 --- sha512crypt (\$6\$)}
Hash no formato \cmd{$6$<sal>$<hash>}:
\begin{lstlisting}[style=output]
$6$aReallyHardSalt$6WKUTqzq.UQQmrm0p/T7MPpMbGNnzXPMAXi4bJMl9be.cfi3/qxIf.hsGpS41BqMhSrHVXgMpdjS6xeKZAs02.
\end{lstlisting}
O prefixo \cmd{$6$} identifica \texttt{sha512crypt} (\cmd{-m 1800}), o esquema do
\cmd{/etc/shadow} moderno; o sal (\cmd{aReallyHardSalt}) já está embutido no
próprio hash, sem necessidade de informá-lo à parte. Por ser um algoritmo com
milhares de iterações (muito lento), o \texttt{rockyou} foi filtrado para
palavras alfanuméricas de \textbf{6 caracteres}; mesmo assim a quebra levou
$\approx 10$ minutos.

\begin{lstlisting}[language=bash,caption={Filtragem por 6 caracteres e quebra do sha512crypt.}]
grep -xE '[A-Za-z0-9]{6}' /usr/share/wordlists/rockyou.txt > words6.txt
hashcat -m 1800 -a 0 t2_3.txt words6.txt -O
\end{lstlisting}
\noindent Texto-claro: \textbf{waka99}.

\subsubsection{\#4 --- HMAC-SHA1 (chave = sal)}
Par \cmd{hash:sal}: \cmd{e5d8870e5bdd26602cab8dbe07a942c8669e56d6:tryhackme}. O
\texttt{hashid} sugere a família SHA-1. A questão é \emph{como} o sal entra: após
descartar as concatenações simples (senha+sal e sal+senha), o hash corresponde a
\textbf{HMAC-SHA1 com o sal como chave} (\cmd{-m 160}).

\begin{lstlisting}[language=bash,caption={HMAC-SHA1 (modo 160) com o sal como chave.}]
hashcat -m 160 -a 0 t2_4.txt /usr/share/wordlists/rockyou.txt \
        -r /usr/share/hashcat/rules/best66.rule
\end{lstlisting}
\noindent Texto-claro: \textbf{481616481616}.

\begin{becosemsaida}
\begin{itemize}[leftmargin=1.4em,itemsep=2pt,topsep=2pt]
  \item \cmd{-m 110} --- \texttt{sha1(senha.sal)}: sem resultado.
  \item \cmd{-m 120} --- \texttt{sha1(sal.senha)}: sem resultado.
  \item \cmd{-m 160} --- \texttt{HMAC-SHA1(chave = sal)}: correspondência.
\end{itemize}
\emph{Lição:} com sal, não basta o algoritmo-base --- é preciso descobrir o
esquema de combinação (concatenação vs.\ HMAC, e a ordem/chave).
\end{becosemsaida}

%--------------------------- HASHES ADICIONAIS --------------------------------
\subsection{Hashes Adicionais (exercício complementar)}

Conjunto extra proposto como prática, aplicando o mesmo fluxo. Alguns vêm com
\textbf{dica de formato}, o que favorece o ataque por \emph{máscara}.

\begin{table}[H]
  \centering
  \caption{Resumo dos hashes adicionais.}
  \label{tab:add}
  \begin{tabular}{l l >{\ttfamily}c l l}
    \hline
    \textbf{Hash (abrev.)} & \textbf{Algoritmo} & \textbf{-m} & \textbf{Ataque} & \textbf{Texto-claro} \\
    \hline
    482c\ldots 1e38        & MD5       & 0    & dicionário  & \texttt{password123} \\
    861c\ldots 5228        & SHA-1     & 100  & dicionário  & \texttt{easy} \\
    4149\ldots 4099        & MD4       & 900  & máscara     & \texttt{0ffs3c} \\
    cdeb\ldots 5875 : sal  & HMAC-SHA1 & 160  & dicionário  & \texttt{ovelha} \\
    362f\ldots 2912        & SHA-256   & 1400 & máscara     & \texttt{sawctf} \\
    \hline
  \end{tabular}
\end{table}

\subsubsection{Add \#1 --- MD5}
\cmd{482c811da5d5b4bc6d497ffa98491e38} (MD5, \cmd{-m 0}).
\begin{lstlisting}[language=bash]
hashcat -m 0 -a 0 add1.txt /usr/share/wordlists/rockyou.txt
\end{lstlisting}
\noindent Texto-claro: \textbf{password123}.

\subsubsection{Add \#2 --- SHA-1}
\cmd{861c4f67e887dec85292d36ab05cd7a1a7275228} (SHA-1, \cmd{-m 100}).
\begin{lstlisting}[language=bash]
hashcat -m 100 -a 0 add2.txt /usr/share/wordlists/rockyou.txt
\end{lstlisting}
\noindent Texto-claro: \textbf{easy}.

\subsubsection{Add \#3 --- MD4 (máscara)}
\cmd{4149c5cc4c378444d116d65ad5ba4099}, com dica ``\emph{apenas letras e números,
6 caracteres}''. Isso pede um ataque de \textbf{máscara} (\cmd{-a 3}) com um
\emph{charset} customizado que una dígitos, maiúsculas e minúsculas
(\cmd{-1 ?d?u?l}). A máscara MD5 esgota sem achar; o algoritmo correto é
\textbf{MD4} (\cmd{-m 900}).
\begin{lstlisting}[language=bash,caption={Máscara de 6 posições (?d/?u/?l) sobre MD4.}]
hashcat -m 900 -a 3 add3.txt -1 ?d?u?l ?1?1?1?1?1?1
\end{lstlisting}
\noindent Texto-claro: \textbf{0ffs3c}.

\begin{becosemsaida}
A mesma máscara sobre MD5 (\cmd{-m 0}) esgotou o espaço sem correspondência.
\end{becosemsaida}

\subsubsection{Add \#4 --- HMAC-SHA1}
Par \cmd{cdeb746ec095149627348b61d4140fc58b745875:satech}, com dica explícita
\textbf{HMAC-SHA1} (\cmd{-m 160}, sal como chave).
\begin{lstlisting}[language=bash]
hashcat -m 160 -a 0 add4.txt /usr/share/wordlists/rockyou.txt
\end{lstlisting}
\noindent Texto-claro: \textbf{ovelha}.

\subsubsection{Add \#5 --- SHA-256 (máscara)}
\cmd{362fda2183b7ac73400a83f6ab2c359451e48adf6c3d46a2963ee2abdf852912} (64 hex,
SHA-256, \cmd{-m 1400}), com dica ``\emph{minúsculas e números, 6 caracteres}''.
O \emph{charset} customizado restringe-se a \cmd{?d?l}, o que torna a máscara
pequena e rápida.
\begin{lstlisting}[language=bash,caption={Máscara restrita a minúsculas e dígitos sobre SHA-256.}]
hashcat -m 1400 -a 3 add5.txt -1 ?d?l ?1?1?1?1?1?1
\end{lstlisting}
\noindent Texto-claro: \textbf{sawctf}.

%============================= 2.7 CONCLUSAO ====================================
\section{Conclusão}
O desafio reforça que a etapa decisiva do \emph{cracking} é a
\textbf{identificação} correta do algoritmo. O segundo fator é a \textbf{escolha da estratégia de ataque}
segundo o custo do algoritmo e o que se sabe da senha: dicionário direto para
palavras comuns; dicionário com \textbf{regras} para mutações
(\emph{Eternity22}, \emph{paule}); \textbf{máscara} quando há dica de formato
(\emph{0ffs3c}, \emph{sawctf}); e, para algoritmos lentos (\texttt{bcrypt},
\texttt{sha512crypt}), \textbf{redução prévia do dicionário} por comprimento para
tornar a quebra viável. Por fim, hashes com \textbf{sal} exigem descobrir o
esquema de combinação (concatenação vs.\ HMAC), como no caso HMAC-SHA1 (modo
160). Por fim, todos os 14 hashes, 5 do Nível 1, 4 do Nível 2 e 5 adicionais, foram
recuperados.

%============================= 2.8 REFERENCIAS ==================================
\section{Referências}
\begin{itemize}[leftmargin=1.5em]
  \item TryHackMe --- Sala \emph{Crack The Hash} ---
        \url{https://tryhackme.com/room/crackthehash}
  \item hashcat --- \emph{example hashes} (identificação de modos) ---
        \url{https://hashcat.net/wiki/doku.php?id=example_hashes}
  \item hashcat --- documentação (Kali Tools) ---
        \url{https://www.kali.org/tools/hashcat/}
  \item CyberChef (função \emph{Magic}) ---
        \url{https://gchq.github.io/CyberChef/}
  \item SecLists --- \emph{Password lists} ---
        \url{https://github.com/danielmiessler/SecLists}
\end{itemize}

\end{document}
